% SHARD-ID: AQM-06-TORIC-CHAIN-GHOST-PROOF
% SHARD-TITLE: Proof of the Toric Chain-to-Ghost Lift
% SHARD-SUMMARY: Gives a concrete theorem relating cellular cohomology, CSS stabilizers, and the ghost supercharge.
% SHARD-SUMMARY: Constructs an explicit cohomology-class basis for the toric-code space.
% SHARD-SUMMARY: Proves nilpotence and the stabilizer Hamiltonian anticommutator from projectors and CAR.
% SHARD-KEYWORDS: toric code, proof, chain complex, CSS code, ghost supercharge, cohomology

\section{Proof of the Toric Chain-to-Ghost Lift}
\label{sec:chain-ghost-proof}

This section replaces the slogan ``the chain complex supplies the supports''
by an explicit finite-dimensional statement.

\subsection{Data}

Fix the \(k\times k\) periodic square cellulation of the torus, \(k\geq 2\).
All chain groups in this subsection are vector spaces over \(\mathbb{F}_2\):
\[
  C_2 \xrightarrow{\partial_2} C_1 \xrightarrow{\partial_1} C_0 .
\]
The bases are faces, edges, and vertices. Thus
\[
  |C_0|=k^2,\qquad |C_1|=2k^2,\qquad |C_2|=k^2 .
\]
Let
\[
  \delta^0=\partial_1^T:C^0\to C^1,\qquad
  \delta^1=\partial_2^T:C^1\to C^2.
\]
The cellular supercharge is the degree \(+1\) map
\[
  Q_{\mathrm{cell}}=\delta^0+\delta^1
  \quad\text{on}\quad C^0\oplus C^1\oplus C^2 .
\]
Since \(\partial_1\partial_2=0\), one has
\[
  Q_{\mathrm{cell}}^2=\delta^1\delta^0=(\partial_1\partial_2)^T=0.
\]
The degree-one cellular cohomology is the quotient
\[
  H^1_{\mathrm{cell}}
  =\ker(\delta^1)/\operatorname{im}(\delta^0).
\]
In elementary terms, its elements are edge-bit strings \(a\) whose parity around
every plaquette is zero, with two such strings identified if they differ by the
edge-bit string produced by applying star flips at some set of vertices.

\subsection{From Boundary Matrices to Pauli Checks}

The physical Hilbert space is
\[
  \mathcal H=(\mathbb C^2)^{\otimes C_1}.
\]
Its computational \(Z\)-basis is written \(\{|a\rangle:a\in C^1\}\), where
\(a_e\in\mathbb{F}_2\) is the bit on edge \(e\).

For an edge-bit string \(r\in C^1\), define Pauli monomials by
\[
  X^r|a\rangle=|a+r\rangle,\qquad
  Z^r|a\rangle=(-1)^{r\cdot a}|a\rangle ,
\]
where \(r\cdot a=\sum_e r_ea_e\) is computed in \(\mathbb{F}_2\).
They obey
\[
  X^rZ^s=(-1)^{r\cdot s}Z^sX^r .
\]

For a vertex \(v\), let \(r_v=\delta^0 v\in C^1\), the row of
\(\partial_1\) supported on edges incident to \(v\). Define the star check
\[
  A_v=X^{r_v}.
\]
For a face \(f\), let \(s_f=\partial_2 f\in C_1\), identified with its edge-bit
string in \(C^1\). Define the plaquette check
\[
  B_f=Z^{s_f}.
\]
Then
\[
  A_vB_f=(-1)^{r_v\cdot s_f}B_fA_v.
\]
But \(r_v\cdot s_f\) is exactly the coefficient of \(v\) in
\(\partial_1\partial_2f\), hence it is zero. Therefore every \(A_v\) commutes
with every \(B_f\). Checks of the same Pauli type commute automatically, and
each check squares to the identity.

\subsection{The Code Space as a Linearized Cohomology Set}

The toric-code space is
\[
  \mathcal C=\{|\psi\rangle\in\mathcal H:
  A_v|\psi\rangle=|\psi\rangle,\ B_f|\psi\rangle=|\psi\rangle
  \text{ for all }v,f\}.
\]
We now construct a basis of \(\mathcal C\) from
\(\ker(\delta^1)/\operatorname{im}(\delta^0)\).

First, for a basis vector \(|a\rangle\),
\[
  B_f|a\rangle=(-1)^{s_f\cdot a}|a\rangle .
\]
The condition \(B_f|a\rangle=|a\rangle\) for every \(f\) is therefore
\[
  s_f\cdot a=0\quad\text{for all }f,
\]
which is precisely \(\delta^1a=0\). Thus the simultaneous \(+1\) eigenspace of
all plaquettes is spanned by \(|a\rangle\) with \(a\in\ker(\delta^1)\).

Second, a star acts by
\[
  A_v|a\rangle=|a+\delta^0 v\rangle.
\]
Thus the star group translates \(a\) by elements of
\(\operatorname{im}(\delta^0)\). Because \(\delta^1\delta^0=0\), these
translations preserve \(\ker(\delta^1)\).

For each class \([a]\in H^1_{\mathrm{cell}}\), define
\[
  |[a]\rangle
  =
  \frac{1}{\sqrt{|\operatorname{im}\delta^0|}}
  \sum_{u\in\operatorname{im}(\delta^0)} |a+u\rangle .
\]
This vector is well-defined: replacing \(a\) by \(a+u_0\) just reindexes the
sum. It is fixed by every star \(A_v\), because \(A_v\) translates the summation
set by \(\delta^0v\). It is fixed by every plaquette \(B_f\), because every
\(a+u\) in the sum lies in \(\ker(\delta^1)\). Hence \(|[a]\rangle\in\mathcal C\).

If \([a]\neq[b]\), the two sums have disjoint computational-basis supports, so
\(|[a]\rangle\) and \(|[b]\rangle\) are orthogonal. Conversely, any vector in
the plaquette \(+1\) eigenspace is a linear combination of basis states with
\(a\in\ker(\delta^1)\), and imposing all star equations forces its coefficients
to be constant on each coset of \(\operatorname{im}(\delta^0)\). Therefore the
vectors \(|[a]\rangle\) form a basis of \(\mathcal C\), and there is a concrete
isomorphism
\[
  \mathbb C[H^1_{\mathrm{cell}}]\cong \mathcal C,\qquad
  [a]\mapsto |[a]\rangle .
\]

For the periodic square torus this gives four states. Indeed, suppose
\(\sum_v \alpha_v(\text{row }v\text{ of }\partial_1)=0\). For each edge with
endpoints \(v,w\), the coefficient of that edge gives
\(\alpha_v+\alpha_w=0\). The graph is connected, so all \(\alpha_v\) are equal.
Thus the only row relation is the sum of all rows, and
\(\operatorname{rank}\partial_1=k^2-1\).

Similarly, suppose \(\sum_f\beta_f\partial_2 f=0\). For each edge separating
adjacent faces \(f,g\), the coefficient of that edge gives
\(\beta_f+\beta_g=0\). The dual graph is connected, so all \(\beta_f\) are
equal. Thus the only column relation is the sum of all face boundaries, and
\(\operatorname{rank}\partial_2=k^2-1\). Hence
\[
  \dim_{\mathbb F_2}H^1_{\mathrm{cell}}
  =2k^2-(k^2-1)-(k^2-1)=2.
\]
Thus \(\dim_{\mathbb C}\mathcal C=|H^1_{\mathrm{cell}}|=2^2=4\).

\subsection{The Ghost Supercharge Is the Koszul Lift of These Checks}

Let the list of stabilizer checks be
\[
  \{S_i\}_{i\in I}=\{A_v:v\in C_0\}\cup\{B_f:f\in C_2\}.
\]
The previous subsection proved that all \(S_i\) commute and that \(S_i^2=I\).
Define the violated-check projectors
\[
  P_i=\frac{I-S_i}{2}.
\]
They are self-adjoint idempotents, and they commute with each other.

Now attach one auxiliary fermion to each \(i\in I\). Let \(\epsilon_i\) create
the \(i\)-ghost and \(\iota_i\) annihilate it. These operations act on the
exterior algebra on the ghost labels and satisfy
\[
  \iota_i\epsilon_j+\epsilon_j\iota_i=\delta_{ij},\qquad
  \epsilon_i\epsilon_j+\epsilon_j\epsilon_i=0,\qquad
  \iota_i\iota_j+\iota_j\iota_i=0 .
\]
They commute with the physical projectors \(P_i\). Define
\[
  Q_{\mathrm{gh}}=\sum_{i\in I}\epsilon_iP_i,\qquad
  Q_{\mathrm{gh}}^*=\sum_{i\in I}\iota_iP_i .
\]

Nilpotence is immediate but worth writing out:
\[
  Q_{\mathrm{gh}}^2
  =\sum_{i,j}\epsilon_i\epsilon_jP_iP_j
  =\sum_{i<j}(\epsilon_i\epsilon_j+\epsilon_j\epsilon_i)P_iP_j=0.
\]
For the anticommutator,
\[
\begin{aligned}
  Q_{\mathrm{gh}}Q_{\mathrm{gh}}^*
  +Q_{\mathrm{gh}}^*Q_{\mathrm{gh}}
  &=
  \sum_{i,j}(\epsilon_i\iota_j+\iota_j\epsilon_i)P_iP_j  \\
  &=\sum_i P_i^2
  =\sum_i P_i .
\end{aligned}
\]
This is the shifted toric-code Hamiltonian on \(\mathcal H\), tensored with the
identity on the ghost exterior algebra.

Finally, degree-zero means ``no ghost present''. A degree-zero vector
\(|\psi\rangle\in\mathcal H\) is \(Q_{\mathrm{gh}}\)-closed exactly when
\[
  Q_{\mathrm{gh}}|\psi\rangle
  =\sum_i \epsilon_iP_i|\psi\rangle=0,
\]
which, by independence of the one-ghost labels, is equivalent to
\(P_i|\psi\rangle=0\) for every \(i\). This is equivalent to
\(S_i|\psi\rangle=|\psi\rangle\) for every star and plaquette. There are no
degree \(-1\) states, so no degree-zero exact states. Therefore
\[
  H^0(Q_{\mathrm{gh}})=\mathcal C
  \cong \mathbb C[H^1_{\mathrm{cell}}].
\]

This is the concrete unification. \(Q_{\mathrm{cell}}\) is the cellular
differential on finite \(\mathbb F_2\)-cochains. \(Q_{\mathrm{gh}}\) is a
complex-linear operator on the physical Hilbert space tensored with ghosts. The
first one supplies the quotient set \(H^1_{\mathrm{cell}}\); the second one has
the toric-code Hamiltonian as its anticommutator and selects the complex vector
space freely spanned by that quotient. The sharper operator-level relation is
derived in Section~\ref{sec:operator-relation}.
